\SourceLeaf{102}
\[
\begin{array}{r|r|r}
A&3A+7B&\mathrm{Fehler}\\\hline
0{,}05&2{,}0869&+0{,}0005\\
0{,}06&2{,}0847&-0{,}0017
\end{array}
\]
Sodann ergeben die fünfzifrigen Tafeln
\[
\begin{array}{r|r|r}
0{,}052&2{,}08667&+0{,}00029\\
0{,}053&2{,}08638&0
\end{array}
\]
Endlich die siebenzifrigen
\[
\begin{array}{r|r|r}
0{,}0529&2{,}0863943&+0{,}0000118\\
0{,}0530&2{,}0863660&-0{,}0000165
\end{array}
\]
Hienach wird
\[
\begin{array}{rcl}
A&=&0{,}0529417\\
\log f&=&2{,}6812412\\
7\log y&=&2{,}7341829\\
\log y&=&0{,}3905976\\
-y&=x=&-2{,}4580892
\end{array}
\]
Für die zweite Auflösung steht die Rechnung, auf ähnliche Weise geführt, folgendermaassen:
\[
\begin{array}{r|r|r}
A&3A+7B&\mathrm{Fehler}\\\hline
0{,}19&2{,}0843&-0{,}0021\\
0{,}20&2{,}0868&+0{,}0004
\end{array}
\]
\[
\begin{array}{r|r|r}
0{,}197&2{,}08627&-0{,}00011\\
0{,}198&2{,}08654&+0{,}00016
\end{array}
\]
\[
\begin{array}{r|r|r}
0{,}1975&2{,}0863805&-0{,}0000020\\
0{,}1976&2{,}0864082&+0{,}0000257
\end{array}
\]
\[
\begin{array}{rcl}
A&=&0{,}1975072\\
\log f&=&2.6812412\\\hline
7\log y&=&2{,}8787484\\
\log y&=&0{,}4112498\\
x&=&-2{,}5778036
\end{array}
\]
\SourceLeaf{103}
Die Gleichung, welche uns hier als Beispiel gedient hat, ist absichtlich so
gewählt, dass zwei ihrer Wurzeln wenig verschieden sind. In einem solchen
Falle sind, wie schon oben im Art.\ 11 bemerkt ist, die Reihen wegen ihrer sehr
langsamen Convergenz wenig brauchbar: auch bei der indirecten Auflösung ist
davon wenigstens eine schwache Analogie erkennbar, indem das Fortschreiten der
successiven Annäherungen bei den beiden negativen Wurzeln (welche eben die
wenig ungleichen sind) etwas träger ist, als bei der positiven. Ein wesentlicher
Unterschied ist aber der, dass die sehr langsame Convergenz der Reihen für sämmtliche Wurzeln eintritt, während bei dem indirecten Verfahren die, auch nur in
geringem Grade fühlbare, langsamere Annäherung lediglich bei den zwei wenig
verschiedenen Wurzeln vorkommt.

\FifthArticle{17}
Ganz verschieden von dem in den vorhergehenden Artikeln gelehrten Verfahren ist dasjenige, welches zur Bestimmung der imaginären Wurzeln angewandt
werden muss. Im Allgemeinen ist die Bestimmung der imaginären Wurzeln auf
indirectem Wege deswegen weit schwieriger, als die der reellen, weil jene aus
einem unendlichen Gebiet von zwei Dimensionen herausgesucht werden müssen,
diese nur aus einem Unendlichen von Einer Dimension, und gerade darum verdient ein sehr umfassender besonderer Fall, wo man jene Schwierigkeit umgehen
und die Frage in dasselbe Gebiet versetzen kann, zu welchem die Aufsuchung
der reellen Wurzeln gehört, eine eigne Ausführung. Einen solchen Fall bieten
die Gleichungen mit drei Gliedern dar.

Da die Methode mit gleicher Leichtigkeit angewandt werden kann, die Coëfficienten der Gleichung mögen reell oder imaginär sein, so lege ich sofort die allgemeine Form der Gleichung zum Grunde
\[
\mathrm{X}=x^{m+n}+e(\cos\varepsilon+i\sin\varepsilon)x^m+f(\cos\varphi+i\sin\varphi)=0
\]
wo $e$ und $f$ positive Grössen bedeuten: für einen reellen, positiven oder negativen, Coëfficienten ist dann der betreffende Winkel ($\varepsilon$ oder $\varphi$) entweder $0$ oder
$180^{0}$. Die Voraussetzung, dass $m$ und $n$ keinen gemeinschaftlichen Divisor haben, wird ohne Beeinträchtigung der Allgemeinheit auch hier beibehalten bleiben können. Eine der Gleichung Genüge leistende imaginäre Wurzel $x=t+iu$
setzt man in die Form $r(\cos\rho+i\sin\rho)$, wobei es für unsern gegenwärtigen Zweck
\SourceLeaf{104}
vortheilhafter ist, die sonst gewöhnliche Bedingung, dass $r$ positiv sein soll, hier
nicht zu machen, sondern anstatt derselben die, dass $\rho$ immer zwischen den Grenzen $0$ und $180^{0}$ genommen werden soll. In dem Fall, wo die Coëfficienten der
Gleichung beide reell sind, kann man den Umfang der Werthe von $\rho$ noch weiter auf die Hälfte verengen: denn da bekanntlich von den imaginären Wurzeln
einer solchen Gleichung je zwei zusammengehören, wie $t+iu$ und $t-iu$, so
wird offenbar für die eine Wurzel jedes Paars der Werth von $\rho$ zwischen $0$ und
$90^{0}$ fallen, und man braucht durch das indirecte Verfahren nur diese zu bestimmen, indem daraus die andere von selbst folgt durch Vertauschung von $\rho$ mit
$180^{0}-\rho$ und von $r$ mit $-r$.

\FifthArticle{18}
Das Wesen der Methode besteht in der Aufstellung einer Gleichung, welche
bloss $\rho$ ohne $r$ enthält. Um dazu zu gelangen, setze man die Gleichung $\mathrm{X}=0$
durch Division mit ihrem ersten Gliede in die Form
\[
1+e(\cos\varepsilon+i\sin\varepsilon)x^{-n}+f(\cos\varphi+i\sin\varphi)x^{-m-n}=0
\]
oder
\[
\begin{aligned}&1+er^{-n}(\cos(n\rho-\varepsilon)-i\sin(n\rho-\varepsilon))\\
&\quad+fr^{-m-n}(\cos[(m+n)\rho-\varphi]-i\sin[(m+n)\rho-\varphi])=0\end{aligned}
\]
Da nun hier die imaginären Theile einander aufheben müssen, so hat man\nobreak\hspace{.5em}(I)
\[
r^m=-\frac{f\sin((m+n)\rho-\varphi)}{e\sin(n\rho-\varepsilon)}
\]
Auf ähnliche Art erhält man, wenn die Gleichung $\mathrm{X}=0$ mit ihrem zweiten
oder dritten Gliede dividirt, und erwägt, dass in beiden Fällen die imaginären
Theile der neuen Gleichungen einander aufheben müssen, die Gleichungen
\[
\left.
\begin{array}{rcl}
r^{m+n}&=&\displaystyle\frac{f\sin(m\rho+\varepsilon-\varphi)}{\sin(n\rho-\varepsilon)}\\[1.2ex]
r^n&=&\displaystyle-\frac{e\sin(m\rho+\varepsilon-\varphi)}{\sin((m+n)\rho-\varphi)}
\end{array}
\right\}\quad\mathrm{(I)}
\]
Man sieht, dass jede der drei Gleichungen (I) auch schon aus der Verbindung der
beiden andern abgeleitet werden kann. Eliminirt man aber $r$ aus Verbindung
zweier, so erhält man (II)
\SourceLeaf{105}
\[
\lambda=(-1)^{m+n}\frac{\sin(m\rho+\varepsilon-\varphi)^m\sin(n\rho-\varepsilon)^n}{\sin((m+n)\rho-\varphi)^{m+n}}
\]
wo zur Abkürzung (eben so wie oben)
\[
\frac{f^n}{e^{m+n}}=\lambda
\]
gesetzt ist. Aus dieser Gleichung hat man die verschiedenen Werthe von $\rho$ zu
bestimmen; den Werth von $r$, welcher jedem Werthe von $\rho$ entspricht, findet
man sodann aus einer der Gleichungen (I), am besten aus der zweiten, rücksichtlich der absoluten Grösse, wobei jedoch in dem Falle, wo $m+n$ gerade ist, noch
eine der beiden andern Gleichungen zur Entscheidung des Zeichens hinzugezogen
werden muss.

\FifthArticle{19}
Die Auflösung der Gleichung II auf indirectem Wege wird man immer mit
Leichtigkeit beschaffen können, wozu noch die Berücksichtigung der folgenden
Bemerkungen beitragen wird.

1) Die Werthe von $\rho$ liegen zwischen $0$ und $180^{0}$; in dem Falle, wo die
Coëfficienten der vorgegebenen Gleichung reell sind, braucht man nur die halbe
Anzahl, nemlich die zwischen $0$ und $90^{0}$ liegenden. einzeln aufzusuchen.

2) In dem einen wie in dem andern Falle wird man zuerst das betreffende
Intervall in die verschiedenen Unterabtheilungen scheiden, die sich durch die
Zeichenabwechslungen in den Werthen der auf der rechten Seite der Gleichung
II stehenden Function von $\rho$ bilden. Die Uebergangswerthe von $\rho$ können offenbar nur solche sein, wo einer der Winkel $m\rho+\varepsilon-\varphi,\ n\rho-\varepsilon,\ (m+n)\rho-\varphi$
durch $180^{0}$ theilbar, und also jene Function selbst entweder $0$ oder unendlich
wird. Von jenen Unterabtheilungen bleiben dann diejenigen, in welchen der
Werth der Function negativ wird, schon von selbst aus der weitern Untersuchung
ausgeschlossen.

3) Falls man nicht schon auf andern Wegen genäherte Werthe von $\rho$ erlangen kann, wird man sich das indirecte Durchsuchen der geeigneten Intervalle
dadurch sehr erleichtern, dass man auf ähnliche Weise, wie aus den Beispielen
des 16. Artikels zu ersehen ist, die ersten Versuche nach abgekürzten Tafeln mit
wenigen Zifern ausführt, und in manchen Fällen möchte man wohl bequem
\SourceLeaf{106}
finden, zuerst nur die Sinuslogarithmen mit drei Zifern auf einem Blättchen etwa
von Grad zu Grad verzeichnet zu diesem Zweck zu verwenden.

\FifthArticle{20}
Zu weiterer Erläuterung mag die Berechnung der imaginären Wurzeln der
oben behandelten Gleichung
\[
x^7+28x^4-480=0
\]
als Beispiel dienen. Nach der Bezeichnung des Art. 17 haben wir hier zuvörderst,
wie oben, $m=4,\ n=3,\ e=28,\ f=480$, und sodann weiter $\varepsilon=0,\ \varphi=180^{0}$.
Die Formeln I des Art. 18 werden demnach
\[
\begin{aligned}
r^4&=\frac{480\sin7\rho}{28\sin3\rho}\\
r^7&=-\frac{480\sin4\rho}{\sin3\rho}\\
r^3&=-\frac{28\sin4\rho}{\sin7\rho}
\end{aligned}
\]
und die Formel II
\[
\frac{1}{\lambda}=\frac{823543}{6750}=\frac{\sin7\rho^7}{\sin3\rho^3\sin4\rho^4}
\]
aus welcher Gleichung zwei zwischen $0$ und $90^{0}$ liegende Werthe von $\rho$ zu bestimmen sind, da die Gleichung $\mathrm{X}=0$ neben ihren drei bereits ermittelten reellen Wurzeln noch zwei Paare zusammengehöriger imaginärer hat. Innerhalb
dieser Grenzen wird $\sin7\rho$ dreimal $=0$, nemlich für $\rho=25\frac57$ Grad, $51\frac37$
Grad und $77\frac17$ Grad, wobei $\sin7\rho^7$ jedesmal sein Zeichen ändert; $\sin3\rho$ wird
einmal $=0$ für $\rho=60^{0}$ gleichfalls mit Zeichenwechsel von $\sin3\rho^3$; endlich
$\sin4\rho$ wird einmal $=0$ für $\rho=45^{0}$, aber ohne Zeichenwechsel für $\sin4\rho^4$.
Erwägt man nun noch, dass der Werth von $\frac{\sin7\rho^7}{\sin3\rho^3\sin4\rho^4}$ für $\rho=0$ dem Grenzwerthe $\frac{7^7}{3^3 4^4}$ gleich zu setzen ist, so wird das Verhalten der Werthe jener Function
in den sechs Unterabtheilungen des Zwischenraumes von $0$ bis $90^{0}$ in folgender
Uebersicht zusammengefasst:
\SourceLeaf{107}
\[
\begin{array}{c|c}
\rho=0\ \text{Grad}&\frac{823543}{6912}\\
&+\\
25\frac57&0\\
&-\\
45&\infty\\
&-\\
51\frac37&0\\
&+\\
60&\infty\,.\\
&-\\
77\frac17&0\\
&+\\
90^{0}&\infty
\end{array}
\]
Man erkennt hieraus, dass sowohl im vierten als im sechsten Zwischenraume nothwendig ein der Formel II Genüge leistender Werth von $\rho$ liegen muss,
und eines Mehrern bedarf es für unsern Zweck nicht, da schon von vorne her fest
steht, dass es nur zwei solche Werthe gibt. Die Gleichung II setze ich in die Form
\[
7\log\sin7\rho-3\log\sin3\rho-4\log\sin4\rho=S=2{,}0863825
\]
Die Auffindung des zwischen $51\frac37$ und $60$ Grad liegenden Werthes durch
allmählige Annäherung vermittelst der Tafeln mit $3,\ 4,\ 5,\ 7$ Zifern zeigt folgendes Schema:
\[
\begin{array}{c|c|c}
\rho&S&\text{Fehler}\\
57^{0}&1{,}527&-0{,}559\\
58&2{,}354&+0{,}268
\end{array}
\]
\[
\begin{array}{c|c|c}
57^{0}40'&2{,}0624&-0{,}0240\\
57\ 50&2{,}2057&+0{,}1193
\end{array}
\]
\[
\begin{array}{c|c|c}
57^{0}41'&2{,}07658&-0{,}00980\\
57\ 42&2{,}09074&+0{,}00436
\end{array}
\]
\[
\begin{array}{c|c|c}
57^{0}41'41''&2{,}0862962&-0{,}0000863\\
57\ 41\ 42&2{,}0865320&+0{,}0001495
\end{array}
\]
Hieraus $\rho=57^{0}41'41''366$, und ferner nach der zweiten der Formeln I,
\SourceLeaf{108}
\[
\begin{array}{lrcl}
\log\sin4\rho&&=&9{,}8891425\,n\\
\mathrm{Compl.}\,\log\sin3\rho&&=&0{,}9193523\\
\log(-480)&&=&2{,}6812412\,n\\\hline
7\log r&&=&3{,}4897360\\
\log r&&=&0{,}4985337
\end{array}
\]
und damit
\[
x=+1{,}6843159+2{,}6637914\,i
\]
so wie die andere dazu gehörige Wurzel
\[
x=+1{,}6843159-2{,}6637914\,i
\]

Der andere zwischen $77\frac{1}{7}$ und $90^{0}$ liegende Werth von
$\rho$ wird durch Anwendung von Tafeln mit drei Decimalen als zwischen
$86^{0}$ und $87^{0}$ liegend erkannt. Die Rechnung in gleicher
Gestalt wie im vorhergehenden Falle steht so:
\[
\begin{array}{c|r|r}
\rho&S&\mathrm{Fehler}\\
86^{0}&1{,}885&-0{,}201\\
87&2{,}533&+0{,}447\\
\end{array}
\]
\[
\begin{array}{c|r|r}
86^{0}10'&1{,}9907&-0{,}0957\\
86\;20&2{,}0946&+0{,}0082\\
\end{array}
\]
\[
\begin{array}{c|r|r}
86\;19&2{,}08409&-0{,}00229\\
86\;20&2{,}09447&+0{,}00809\\
\end{array}
\]
\[
\begin{array}{c|r|r}
86^{0}19'13''&2{,}0863229&-0{,}0000596\\
86\;19\;14&2{,}0864970&+0{,}0001145
\end{array}
\]
\[
\rho=86^{0}19'13''342
\]
\[
\begin{array}{lrcl}
\log\sin4\rho&&=&9{,}4049540\,n\\
\mathrm{Compl.}\,\log\sin3\rho&&=&0{,}0081108\,n\\
\log(-480)&&=&2{,}6812412\,n\\\hline
7\log r&&=&2{,}0943060\,n\\
\log r&&=&0{,}2991866\,n
\end{array}
\]
Zieht man vor, $r$ positiv zu haben, so braucht man nur zugleich für
$\rho$ den um $180^{0}$ vergrösserten Werth
$266^{0}19'13''342$ anzusetzen. Die Wurzel selbst ist
\[
x=-0{,}1278113-1{,}9874234\,i
\]
und die andere dazu gehörige nur im Zeichen des imaginären Theils
davon verschieden.
\par\SourceLeaf{109}
Die sämmtlichen Wurzeln der Gleichung $x^7+28x^4-480=0$ sind demnach
\[
\begin{array}{l}
+1{,}9228841\\
-2{,}4580892\\
-2{,}5778036\\
+1{,}6843159+2{,}6637914\,i\\
+1{,}6843159-2{,}6637914\,i\\
-0{,}1278113+1{,}9874234\,i\\
-0{,}1278113-1{,}9874234\,i
\end{array}
\]
Die Summe der Wurzeln $+0{,}0000005$ ist so genau mit dem wahren
Werthe $0$ übereinstimmend, wie nur von dem Gebrauch siebenzifriger
Logarithmen erwartet werden durfte. In der andern Form hat man
\[
\begin{array}{r|r}
\log r&\rho\\\hline
0{,}2839531&0\\
0{,}3905976&180^{0}\\
0{,}4112498&180\\
0{,}4985337&57\;41'41''366\\
0{,}4985337&302\;18\;18{,}634\\
0{,}2991866&93\;40\;46{,}658\\
0{,}2991866&266\;19\;13{,}342
\end{array}
\]
Die Summe der Logarithmen der Werthe von $r$ findet sich
$=2{,}6812411$, gleichfalls befriedigend genau mit dem Logarithmen
von $480$ übereinstimmend.

Es wird übrigens kaum nöthig sein zu erinnern, dass die in diesem so
wie die im 16. Artikel aufgestellten Rechnungen nur dazu bestimmt
sind, den Gang der Arbeit nach ihren Hauptmomenten zu erläutern,
keinesweges aber für die Form des kleinen Mechanismus der Operationen
maassgebend sein sollen. Geübtere Rechner werden meistens vorziehen,
nicht so viele Zwischenstufen anzuwenden, als in jenen Beispielen
geschehen ist. Ueberhaupt wird jeder in dergleichen Arbeiten
einigermassen erfahrene die Einzelheiten des Geschäfts leicht selbst
in diejenige Gestalt bringen, die den jedesmaligen Umständen und
seiner eignen individuellen Gewöhnung am meisten angemessen ist,
und es kann hier nicht der Ort sein, in solche Einzelheiten weiter
einzugehen.
\begin{center}\rule{0.25\linewidth}{0.3pt}\end{center}
